
\section{Nudging: Theory and Application}

\subsection{Theoretical Foundations}

Nudging, also known as Newtonian relaxation or the dynamic initialization technique, is a simple data assimilation (DA) method where a feedback term is added to the dynamical model equations. This feedback term drives the model solution toward the observations.

\subsubsection*{Continuous Time Formulation}

Assume the state variable $\mathbf{x}(t) \in \mathbb{R}^n$ satisfies:
\[
\frac{d\mathbf{x}}{dt} = \mathcal{M}(\mathbf{x}(t)), \quad \mathbf{x}(0) = \mathbf{x}_0
\]
We observe the system as $\mathbf{y}^{o}(t) \in \mathbb{R}^p$ and define the observation operator $\mathcal{H} : \mathbb{R}^n \to \mathbb{R}^p$, so the observation error is:
\[
\boldsymbol{\epsilon}^o(t) = \mathcal{H}(\mathbf{x}(t)) - \mathbf{y}^{o}(t)
\]
Nudging adds a corrective term to obtain:
\[
\frac{d\mathbf{x}}{dt} = \mathcal{M}(\mathbf{x}(t)) - \mathbf{K}(\mathcal{H}(\mathbf{x}(t)) - \mathbf{y}^{o}(t))
\]
$\mathbf{K}$ is a gain matrix that controls the relaxation strength.

\subsubsection*{Linear Case}

Let $\mathcal{M} = \mathbf{M}$ (a matrix), and $\mathcal{H} = \mathbf{I}$:
\[
\frac{d\mathbf{x}}{dt} = \mathbf{M} \mathbf{x} - \mathbf{K} (\mathbf{x} - \mathbf{y}^o)
\]
Its solution is:
\[
\mathbf{x}(t) = e^{-(\mathbf{K} - \mathbf{M})t} \int_0^t e^{(\mathbf{K} - \mathbf{M})s} \mathbf{K} \mathbf{y}^o(s) ds + e^{-(\mathbf{K} - \mathbf{M})t} \mathbf{x}_0
\]

If $\mathbf{K}$ is symmetric positive-definite, then:
\[
\lim_{K \to \infty} \mathbf{x}(t) = \mathbf{y}^o(t)
\]

\subsubsection*{Discrete Time and Kalman Filter Relation}

In discrete time, nudging takes the form:
\[
\mathbf{x}_{k+1} = \mathbf{M}_{k+1:k} \mathbf{x}_k - \mathbf{K} \left[ \mathcal{H}(\mathbf{M}_{k+1:k}(\mathbf{x}_k)) - \mathbf{y}_{k+1}^o \right]
\]
This resembles the Kalman Filter update:
\[
\mathbf{x}_{k+1}^a = \mathbf{x}_{k+1}^f + \mathbf{K} \left[ \mathbf{y}_{k+1}^o - \mathcal{H}(\mathbf{x}_{k+1}^f) \right]
\]

\subsection{Variants and Practical Considerations}

\begin{itemize}
  \item \textbf{Spectral Nudging}: Only large-scale components (low-wavenumber modes) are nudged, allowing small-scale dynamics to evolve freely.
  \item \textbf{Back-and-Forth Nudging (BFN)}: Nudging is performed forward and backward in time, improving convergence over finite assimilation windows.
  \item \textbf{Optimal Nudging}: K is optimized using a variational approach to minimize a cost function involving observation misfit and regularization.
\end{itemize}

\subsection{Application Example: Turbulent Flow Around a Square Cylinder}

Zauner et al. (2020) apply nudging to a URANS simulation of flow around a square cylinder at $Re = 22{,}000$. DNS data is used as observation input to guide the URANS solution.

\subsubsection*{Model Setup}

\begin{itemize}
  \item Governing equations: URANS with $k$--$\omega$ SST model.
  \item Geometry: 2D square cylinder in cross-flow.
  \item Observations: velocity fields from DNS.
  \item Nudging applied: in a spatial region downstream of the cylinder, with $\alpha = 1$.
\end{itemize}

\subsubsection*{Equation}

Nudged momentum equation:
\[
\frac{\partial \mathbf{u}}{\partial t} + (\mathbf{u} \cdot \nabla) \mathbf{u} = -\nabla p + \nu \nabla^2 \mathbf{u} + \alpha (\mathbf{u}_{\text{DNS}} - \mathbf{u})
\]

\subsubsection*{Results}

\begin{itemize}
  \item Vortex shedding frequency ($St \approx 0.137$) was accurately recovered.
  \item High-frequency KH instabilities ($St \approx 4.384$) were also captured.
  \item SPOD analysis showed spectral agreement between nudged URANS and DNS.
\end{itemize}

\textbf{Conclusion:} Nudging allowed URANS to capture multiscale flow features with minimal complexity, bridging the gap between low-cost and high-fidelity simulations.

